\documentclass[a4paper,12pt]{article}
\usepackage[utf8]{inputenc}
\usepackage[T1]{fontenc}
\usepackage[ngerman]{babel}
\usepackage{amsmath, amssymb}
\usepackage{geometry}
\geometry{left=3cm, right=3cm, top=2.5cm, bottom=2.5cm}
\usepackage{array, booktabs}
\usepackage{enumitem}
\usepackage{parskip}

% Glossar-Eintrag: {Nummer}{Fetter Titel}{Inhalt}
% Für Einträge ohne eigene Nummer: \gentry{}{Titel}{Inhalt}
\newcommand{\gentry}[3]{%
  \noindent\textbf{#1\quad #2}\nopagebreak\\[0.2em]
  #3\par\smallskip
}

\begin{document}

\begin{center}
  {\LARGE\bfseries Glossar zu Lektion N}\\[0.5em]
  {\large MODULNAME (Modul NUMMER)}
\end{center}

\vspace{1em}

Dieses Glossar fasst die wesentlichen Definitionen, Merkregeln und Sätze
der Lektion kompakt zusammen. Nummerierungen entsprechen dem Lehrtext.

% ======================================================================
\section*{1.1 \quad ABSCHNITTSTITEL}

\gentry{1.1.1}{TITEL}{%
  % Definition / Formel / Wahrheitstafel
}

% Wahrheitstafel-Beispiel:
% \gentry{1.2.3}{Negation $\lnot A$}{%
%   \[
%     \begin{array}{c|c}
%       A & \lnot A \\\hline
%       w & f \\ f & w
%     \end{array}
%   \]
% }

% Aufzählung innerhalb eines Eintrags:
% \gentry{}{TITEL}{%
%   \begin{itemize}[topsep=2pt,itemsep=1pt]
%     \item ...
%     \item ...
%   \end{itemize}
% }

% ======================================================================
\section*{1.2 \quad ABSCHNITTSTITEL}

% Weitere Einträge...

\end{document}
